\begin{table*}[t]
\centering\footnotesize
\caption{Retained AUF expert assessment from the original survey, using criterion-specific scores 0--3. These are specification assessments, not measured agent performance or independently validated rankings. U1 identity; U2 types; U3 units; U4 relations; U5 behaviour; U6 constraints; U7 provenance; U8 query; U9 economy; U10 ecosystem.}
\label{tab:auf-matrix}
\setlength{\tabcolsep}{4pt}
\begin{tabular}{@{}llcccccccccc@{}}
\toprule
Layer & Technology & U1 & U2 & U3 & U4 & U5 & U6 & U7 & U8 & U9 & U10 \\
\midrule
L0 & ECLASS / IEC CDD / IEC 61987 & 3&2&2&1&0&1&0&1&2&3 \\
L0 & GS1 identifiers / EPCIS & 3&1&0&1&1&0&2&1&2&3 \\
L1 & OPC UA core and companions & 3&3&2&2&3&1&2&1&1&3 \\
L1 & Asset Administration Shell & 3&3&2&2&2&1&2&1&1&3 \\
L1 & AutomationML & 2&3&1&3&1&1&1&0&1&2 \\
L1 & MTP & 2&3&2&2&3&2&1&0&2&2 \\
L1 & PackML & 1&2&0&0&3&2&0&0&3&3 \\
L1 & ISA-95 / IEC 62264 & 2&2&1&3&1&0&1&0&2&3 \\
L1 & W3C WoT TD 1.1 & 3&3&2&1&3&1&1&1&2&2 \\
L1 & NGSI-LD & 3&2&2&3&1&1&3&2&2&2 \\
L1 & IEC 61850 & 3&3&2&2&2&1&2&0&1&3 \\
L2 & RDF / RDFS & 3&2&0&3&0&0&1&3&1&3 \\
L2 & OWL 2 & 3&3&0&3&0&2&0&2&1&3 \\
L2 & SHACL & 2&3&1&2&1&3&1&2&2&3 \\
L2 & SPARQL 1.1 & 2&1&0&3&0&1&1&3&3&3 \\
L2 & PROV-O & 3&2&0&3&1&0&3&2&2&2 \\
L3 & QUDT & 3&2&3&1&0&2&0&2&2&2 \\
L3 & SOSA / SSN & 3&2&1&2&1&0&3&2&2&3 \\
L3 & Brick & 3&3&1&3&0&1&0&3&2&3 \\
L3 & SAREF and extensions & 3&2&2&2&2&0&1&2&2&2 \\
L3 & DEXPI & 2&3&2&3&0&1&1&1&1&2 \\
L4 & IDS-RAM / Dataspace Protocol & 3&2&0&2&1&2&3&1&1&2 \\
L4 & Catena-X / Manufacturing-X & 3&3&1&2&1&2&3&1&1&2 \\
L4 & ODRL & 2&2&0&2&2&3&2&1&2&2 \\
\bottomrule
\end{tabular}
\end{table*}
\subsection{Assessment rubric and operational failure mapping}
For U1 the key failure is selecting a different physical entity. U2 concerns
malformed calls, invented fields, illegal enumeration members and read-only writes.
U3 separates dimensional incompatibility from compatible units at the wrong scale.
U4 supports containment-based scope and topological impact; U5 describes legal
commands and transitions; U6 publishes checkable guards. U7 binds answers and
decisions to source, time and quality. U8 moves aggregation to a query engine,
U9 asks whether relevant evidence fits the available budget, and U10 records
whether tools, converters and governance permit deployment. A capability can be
valuable even when implemented procedurally: AUF assesses what is made available,
not which syntax deserves credit.

Table~\ref{tab:auf-matrix} retains the original 24 technology assessments. The
criterion-specific ordinal rubric distinguishes absent, partial/indirect,
explicit/incomplete and strongest specified support. Scores are not interval
measurements and must not be averaged into an accuracy prediction. Technology
versions, extensions and exposed profiles can change a judgement. The matrix
is a navigation aid, not a claim that every deployment satisfies all scored
criteria. Type-system support is widespread, whereas unit calculus and executable
constraint coverage concentrate in fewer rows. This motivates composition,
not a universal requirement to adopt every listed technology.

\subsection{An evidence contract for each proposal}
Let $F$ denote available engineering facts and $R(q)$ the declared evidence
requirements of task $q$. The reported coverage and sufficiency are
\begin{equation}
 c(q,F)=\frac{|R(q)\cap F|}{|R(q)|},\quad
 s(q,F)=\mathbf{1}[R(q)\subseteq F].
\end{equation}
These quantify a particular evidence contract; another valid derivation may
need a different requirement set. For an action, validation has three dispositions:
reject if a violation is proved, abstain if required evidence is unresolved, and
accept only relative to the checked contract if neither condition applies.
Acceptance is not a proof that all physical hazards have been modelled.

For compatible affine units, a proposed source value $x$ is converted by
\begin{equation}
 x_t=\frac{m_sx+b_s-b_t}{m_t},
\end{equation}
with dimension compatibility checked first. Engineering bounds apply to $x_t$,
not the source numeral. A faithful numeric answer additionally compares $x_t$
with a referenced observation under a declared tolerance. A writable point's
guard predicates and operating state are separate from that numeric check.

\subsection{Runtime insertion points and engineering ownership}
Semantics enters retrieval, tool-schema synthesis, constrained decoding and
post-proposal verification. Retrieval resolves equipment before projecting
relevant values, relationships and guards. Tool schemas can constrain identifiers,
enumerations and types; only decoder-enforced features eliminate calls at
generation time. Snapshot-dependent conditions still require live validation.
A repair report should name the violated rule and missing evidence; retries
remain bounded and authorization unchanged.

The source model stays behind a query interface. A parameterised service can
answer a plant-wide aggregate, expose current admissible commands or return
an evidence slice with provenance. MCP is one possible transport for such tools,
not a source of equipment truth. Integration must authenticate callers, constrain
scope, preserve mapping versions and avoid unrestricted graph or PLC access.
Figure~\ref{fig:runtime} separates supervisory proposal generation from
controller-side execution authority.

\begin{figure*}[t]
\centering
\begin{tikzpicture}[font=\footnotesize, box/.style={draw,align=center,minimum height=11mm,text width=32mm},arr/.style={-{Stealth},thick}]
\node[box] (source) at (0,0) {Engineering sources\\OPC UA/AAS, recipes, guards};
\node[box] (evidence) at (4,0) {Evidence service\\units, state, time, provenance};
\node[box] (agent) at (8,0) {Agent proposal\\typed tools and bounded repair};
\node[box] (gate) at (12,0) {Dispatch gate\\refresh, scope, approval};
\node[box] (plc) at (12,-1.8) {S7-1500 controller\\current guards and acknowledgement};
\node[box] (rig) at (8,-1.8) {Physical process\\sensors and bounded actuators};
\node[box] (audit) at (4,-1.8) {Independent trace\\proposal, decisions, outcome};
\draw[arr] (source)--(evidence);
\draw[arr] (evidence)--(agent);
\draw[arr] (agent)--(gate);
\draw[arr] (gate)--(plc);
\draw[arr] (plc)--(rig);
\draw[arr] (rig)--(audit);
\draw[arr] (audit)--(evidence);
\end{tikzpicture}
\caption{Proposed supervisory architecture. Evidence supports the agent; dispatch and PLC checks retain current-condition authority. This is a design, not evidence of a completed physical trial.}
\label{fig:runtime}
\end{figure*}